% AKS: new continuation, 14 September 2026.
% This body is self-contained at the level of its explicitly stated inherited inputs.
\section{The absolute common kernel: the new finite result}
\label{sec:AKS}

This complete incoming calculation sharpens the current absolute common
kernel estimate. Its original comparison was against the preceding
logarithmic bound; the intervening KAF chapter already established
the $O(kq)$ scale. The new finite envelope reduces its uniform limsup
constant from $16(4+10\log2)$ to $32\log(25\log3/(4\pi))$.
AKJ below takes the pointwise minimum of both finite bounds. No
limiting kernel coefficient is inferred from these upper bounds.

Retain the actual hypothetical complete quartet and its entire quotient
\[
 \rho=\tfrac12+\delta+i\gamma,\qquad 0<\delta<\tfrac12,\quad \gamma>2,
 \qquad h(z)=\prod_{\alpha\in\{\rho,\bar\rho,1-\rho,1-\bar\rho\}}
 (z-\alpha)^m,\quad v_h=2\xi/h.
\]
The full primary unit is $U=M_{j_hv_h}$. The original tensor order, centre,
polynomial, and coefficient quotient are
\[
\begin{gathered}
 k=4l+1,\quad l\ge1,\quad e=1+k(m-1),\quad q=e(k+1)^2,\quad c=k/2,\\
 \lambda_{a,b}=c+(2a-k)\delta+i(2b-k)\gamma,\qquad
 \chi(S)=\prod_{a,b=0}^k(S-\lambda_{a,b})^e,
 \quad E=\mathbb C[S]/(\chi).
\end{gathered}\tag{AKS1}
\]
For the principal conclusion $m=1$. The source-operator estimates below
also hold with the stated repeated-root list for $m>1$, but no simple-quartet
rank is assigned to those multiplicities.

Let $\Lambda_k(u):E\twoheadrightarrow B_k(u)$ be the original observation,
$K(u)=\ker\Lambda_k(u)$, and $I:K(u)\hookrightarrow E$ its original
coefficient inclusion. For a source $s\in\{1,k\}$, let $G_N^{(s)}$ be the
original minimum-remainder Gram, in the frame $(1,S,\ldots,S^{q-1})$.
Define $H_{K,N}^{(s)}=I^*G_N^{(s)}I$ and retain exactly
\[
 F_K^{(s)}=
 \log\det H_{K,q-1}^{(s)}+\log\det H_{K,q}^{(s)}
 -\log\det H_{K,2q-1}^{(s)}-\log\det H_{K,2q}^{(s)}.
 \tag{AKS2}
\]
Every empty determinant is one. All conjugate transposes refer to the
specified coordinate matrices, not to a bilinear residue pairing.

Here are entirely explicit positive numbers used in the bound:
\[
\begin{gathered}
 R=k\sqrt{\delta^2+\gamma^2},\quad b_s=s/2,\quad
 t_q=\frac\pi2-\frac1q,\quad \mathfrak a_q=\frac{2\log3}{t_q},\quad
 \mathfrak s_q=\frac{\mathfrak a_q^q-1}{\mathfrak a_q-1},\\
 \mathfrak b_{s,q}=\left[\frac9{25}\sin\frac1q\right]^{-b_s}
       \exp\left(2R\left(\log3+\frac{t_q}{2}\right)\right),\\
 \mathcal Z_{s,j}=1+\mathfrak b_{s,q}\mathfrak s_q^2
        \sum_{d=q}^{q+j-1}\frac{d!}{(b_s)_d}
                          \left(\frac{25}{16}\right)^d,
 \quad \mathcal L_{s,j}=\log \mathcal Z_{s,j}\quad (0\le j\le q+1).
\end{gathered}\tag{AKS3}
\]
The empty sum gives $\mathcal Z_{s,0}=1$. In the admitted family $q\ge36$, so
$t_q\in(0,\pi/2)$ and $\mathfrak a_q>1$. These parameters are radii in the existing
Laplace and generating functions; neither is a new source order.

\paragraph{Main theorem (absolute finite bound).}
For the original simple quartet, every original logarithm branch, and
every $|u|\ge R_*$, where $R_*$ is precisely RPC21 or RPC23 (recorded in
Section~\ref{sec:AKS-radius}), put $r=8k-16$. Then
\[
 \boxed{\quad
 0\le F_K^{(s)}(k,u)\le r\bigl(\mathcal L_{s,q}+\mathcal L_{s,q+1}\bigr),
 \qquad s\in\{1,k\}.\quad}
 \tag{AKS4}
\]
There is a stronger, fully period-dependent interval in AKS28 below.
In particular, with
\[
 \kappa_* =32\log\left(\frac{25\log3}{4\pi}\right),
 \qquad 25.02078097649<\kappa_*<25.02078097650,
 \tag{AKS5}
\]
one has the uniform statement
\[
 \boxed{\quad
 0\le\limsup_{k\to\infty\atop k\equiv1\ (4)}
       \ \sup_{\substack{|u|\ge R_*\\\text{original branches}}}
         \frac{F_K^{(s)}(k,u)}{kq}\le\kappa_*.
 \quad}\tag{AKS6}
\]
Here the source may be $s=1$ for every $k$, or $s=k$. The same bound holds
for the actual arithmetic kernel. Its exact finite allowances are in
AKS34--36. Thus the absolute kernel is $O_h(kq)$ on this domain, not only
$O_h(kq\log q)$. Neither AKS5 nor AKS6 calls $\kappa_*$ the actual leading
coefficient.

\section{An exact coefficient realization of the actual lost directions}
\label{sec:AKS-frame}

\subsection{The period matrix and invariant restriction, with their full coefficients}
For $m=1$, order the one-factor roots as $(++),(+-),(-+),(--)$, and let
$\mathcal E_{\rm prim}$ have their original Lagrange idempotents as columns.
Set $\zeta=\exp(2\pi i/5)$ and $F_{jt}=\zeta^{jt}-1$, $1\le j,t\le4$.
The exact inverse is $(F^{-1})_{tj}=\zeta^{-tj}/5$. Retain
\[
\begin{gathered}
 \Pi[P]_j=\int_{\Gamma_j}e^{\Phi_h(z)/u}\operatorname{rem}_hP(z)\,dz,
 \qquad\Phi_h'=h,\quad\Phi_h(0)=0,\\
 \Gamma_j=-\ell_0+\ell_j,\quad
 \arg\ell_j=(\arg u+\pi+2\pi j)/5,\\
 \mathsf H=F^{-1}\Pi U\mathcal E_{\rm prim},\quad
 D=\operatorname{diag}(\zeta,\zeta^2,\zeta^3,\zeta^4),\\
 Q(x)=Q_0(\mathsf H^{-1}x),\qquad
 Q_0(t)=t_1t_4-t_2t_3,\quad Q_a(x)=Q(D^{-a}x).
\end{gathered}\tag{AKS7}
\]
The integral always uses the degree-less-than-four representative.
Invertibility of $\Pi$, all contour phases, and its full determinant are
the PDE/PQO results in the supplied proofs. The four unit values are
$(a,\bar a,\bar a,a)$, with the actual $a=v_h(\rho)\ne0$; no value of $a$
is selected. Hence $\mathsf H$ is the original invertible matrix.

The following formulas make the coefficient realization directly
computable. Write
\[
 \mathcal I_{k,0}=\{\nu\in\mathbb Z_{\ge0}^4:|\nu|=k,
                         \ \nu_1+2\nu_2+3\nu_3+4\nu_4\equiv0\pmod5\}.
\]
For $\nu\in\mathcal I_{k,0}$ define $B_{\nu;(a,b)}$ by summing over
$4\times4$ arrays $n_{t,\alpha}\ge0$ with row sums $\nu_t$, total plus-real
count $a$, and total plus-imaginary count $b$:
\[
 B_{\nu;(a,b)}=
 \sum_{\substack{\sum_\alpha n_{t,\alpha}=\nu_t\\
       \sum_{t,\alpha:\epsilon_\alpha=+}n_{t,\alpha}=a\\
       \sum_{t,\alpha:\eta_\alpha=+}n_{t,\alpha}=b}}
 \left(\prod_{t=1}^4\frac{\nu_t!}{\prod_\alpha n_{t,\alpha}!}\right)
                  \prod_{t,\alpha}\mathsf H_{t,\alpha}^{n_{t,\alpha}}.
 \tag{AKS8}
\]
Let $\mathsf V_{(a,b),d}=\lambda_{a,b}^d$ for $0\le d<q$ and set
$\mathsf A=B\mathsf V$. All entries of $\mathsf A$ are now finite
expressions in the actual periods and unit. There is no generic period
matrix in this definition.

For completeness, put $t_\alpha(w)=\exp((\eta_\alpha\gamma-
 i\epsilon_\alpha\delta)w)$ and $x(w)=\mathsf Ht(w)$. Expansion gives
\[
 x(w)^\nu=\sum_{a,b}B_{\nu;(a,b)}e^{\omega_{a,b}w},\qquad
 \omega_{a,b}=(2b-k)\gamma-i(2a-k)\delta.
\]
In the original primary algebra,
$e^{wY}[1]=\sum_{a,b}e^{\omega_{a,b}w}\varepsilon_{a,b}$, where
$Y=(M_S-cI)/i$. The exponentials are independent by the Vandermonde of
the distinct $\omega_{a,b}$. OCS9's unscaled symmetric word basis
therefore identifies the original observation with
\[
       \jmath\Lambda_k=\mathcal F_0\mathsf A,
       \qquad\ker\mathsf A=K(u)
       \quad\text{in the original coefficient frame.}
       \tag{AKS9}
\]
Every multinomial in AKS8 belongs to its actual ordered words; the
word multiplicity is not inserted a second time. The exact inherited
row Gram of $\mathcal F_0$ is the restriction to $\mathcal I_{k,0}$ of
\[
 \mathcal G_{\nu,\mu}=
 \sum_{\substack{n_{tu}\ge0\\\sum_u n_{tu}=\nu_t\\
                              \sum_t n_{tu}=\mu_u}}
 \frac{k!}{\prod_{t,u}n_{tu}!}\prod_{t,u}
          [5(\delta_{tu}+1)]^{n_{tu}}.
 \tag{AKS10}
\]
This follows by taking the coefficient of $\bar x^\nu y^\mu$ in
$(\bar x^TF^*Fy)^k$, since $F^*F=5(I+\mathbf1\mathbf1^*)$. It is positive
definite and includes every row cross term.

Here is the precise role of the orbit ideal. The map sending a homogeneous
polynomial $f$ to the covector specified by
$\Theta_k(f)(e^{wY}[1])=f(x(w))$ has kernel
$Q\mathbb C[x]_{k-2}$ and is onto $E^\vee$. Indeed equal exponential
indices correspond exactly to monomial differences divisible by
$t_1t_4-t_2t_3$; the displayed Vandermonde proves that there are no other
relations. Its invariant image is $\operatorname{im}\Lambda_k^\vee$.
On the five-member domain, the five $Q_a$ are distinct prime quadrics.
An invariant polynomial divisible by $Q$ is divisible by each $Q_a$,
and therefore by their full product; its quotient is invariant by
cancellation in the polynomial ring. Consequently
\[
\begin{gathered}
 \mathbb C[x]^{\langle D\rangle}\cap (Q)
       =\left(\prod_{a=0}^4Q_a\right)\mathbb C[x]^{\langle D\rangle},\\
 K(u)^\vee\simeq
 \mathbb C[x]_k\big/(\mathbb C[x]_k^{\langle D\rangle}
                              +Q\mathbb C[x]_{k-2}),\\
 b:=\operatorname{rank}\mathsf A
     =d_{5,k,0}-d_{5,k-10,0}=k^2-6k+17,\qquad r=q-b=8k-16.
\end{gathered}\tag{AKS11}
\]
Negative degrees are zero. The count is
$d_{5,d,0}=(\binom{d+3}{3}+4\epsilon_d)/5$,
$\epsilon_d=\mathbf1_{d\equiv0(5)}-\mathbf1_{d\equiv1(5)}$;
subtracting the cubics gives the stated expression for $k\ge10$, and
$k=5,9$ give $12,44$ directly. This is the inherited OQX calculation,
repeated to specify which kernel AKS9 realizes. The source metric is
still to be applied to this kernel, not to the invariant ring itself.

\subsection{A determinant formula requiring no assumed nonzero pivot}
Let $\mathcal G_0$ be the full positive row Gram in AKS10 and set
\[
 W=\mathsf A^*\mathcal G_0\mathsf A,\qquad
 \sigma_j=[z^j]\det(I_q+zW),\qquad
 P_K=\frac1{\sigma_b}\sum_{j=0}^b(-1)^j\sigma_{b-j}W^j.
 \tag{AKS12}
\]
Since $W$ has exactly $b$ positive eigenvalues, $\sigma_b>0$. The scalar
polynomial in AKS12 is one at zero and zero at every nonzero eigenvalue:
it is the normalized product of the factors $\lambda_i-z$ for the $b$
positive eigenvalues. Hence $P_K=P_K^*=P_K^2$ and
$\operatorname{im}P_K=K(u)$. This is the projector derived from the
literal observation, not a freely chosen projection. Replacing
$\mathcal G_0$ by any other positive coordinate Gram leaves this same
coefficient-orthogonal kernel projector, since its image and orthogonal
complement are unchanged. No such replacement is needed here.

At every original endpoint define
\[
 D_{K,N}^{(s)}=[z^r]\det(I_q+zP_KG_N^{(s)}P_K).
 \tag{AKS13}
\]
These numbers are strictly positive for $r>0$ and are one for $r=0$.
To prove their exact metric meaning, choose an orthonormal coefficient
basis $U_K$ of $\operatorname{im}P_K$ just for the proof. In a unitary
extension, $P_KG_NP_K$ has its positive block $U_K^*G_NU_K$ and a zero
block. Thus $D_{K,N}=\det(U_K^*G_NU_K)$. Write the original inclusion
as $I=U_KT$; then $|\det T|^2=\det(I^*I)$ and
\[
 \det H_{K,N}^{(s)}=\det(I^*I)D_{K,N}^{(s)},\qquad
 \boxed{\ F_K^{(s)}=
 \log\frac{D_{K,q-1}^{(s)}D_{K,q}^{(s)}}
               {D_{K,2q-1}^{(s)}D_{K,2q}^{(s)}}.\ }
 \tag{AKS14}
\]
This supplies an equivalent exact coefficient determinant formula for
all the $8k-16$ lost directions, including the entire period dependence.
The original frame factor is present at each endpoint and cancels only
with its four actual signs. It is unnecessary to suppose any particular
minor nonzero. Alternatively, selecting the first independent $r$ columns
of $P_K$ gives an explicit coefficient frame; the determinant formula
avoids even that chart choice.

For $m>1$ replace AKS8--9 by OCS12--18's full confluent matrix
$\mathfrak A_{\nu;(a,b,D)}=i^DD!c_{\nu;a,b,D}$ and its primary isomorphism.
AKS12--14 are then identical with the actual rank, not with the rank in
AKS11. All primary nilpotents and all unit derivatives remain in that
matrix.

\section{A bounded exact remainder operator on the original source}
\label{sec:AKS-operator}

\subsection{The complete source norms and coefficient map}
For either existing source set $M_s=(2\pi)^{s/2}$ and retain
\[
 dm_{0,s}(y)=\frac{M_s2^{b_s}}{4\pi\Gamma(b_s)}
           |\Gamma(b_s/2+iy/2)|^2\,dy,\qquad
 \int e^{ty}dm_{0,s}=M_s(\cos t)^{-b_s},\quad |t|<\pi/2.
 \tag{AKS15}
\]
The supplied IGO1--3 proof derives this transform by the beta integral,
Fourier inversion and a contour shift within $|\operatorname{Im}x|<\pi/2$.
In particular it includes the exact masses $M_1=\sqrt{2\pi}$ and
$M_k=(2\pi)^{k/2}$ and exponential domination on every smaller strip.
The original real monic polynomials satisfy
\[
\begin{gathered}
 p_0=1,\quad p_1=y,\quad
 p_{d+1}=yp_d-d(b_s+d-1)p_{d-1},\\
 \sum_{d\ge0}\frac{p_d(y)}{d!}z^d
      =(1+z^2)^{-b_s/2}e^{y\arctan z},\qquad
 h_d^{(s)}=\|p_d\|^2=M_s d!(b_s)_d.
\end{gathered}\tag{AKS16}
\]
The generating formula follows by differentiating in $z$. To check the
normalization directly, integrate the product of two generating
functions at small real $z,w$. AKS15 and the cosine addition identity
give $M_s(1-zw)^{-b_s}$. Taking any finite pair of derivatives at zero,
justified by the stated exponential domination, gives orthogonality and
the norms in AKS16. Thus the leading phase in
$u_d(S)=p_d((S-c)/i)/\sqrt{h_d^{(s)}}$ is exactly $i^{-d}$.
The same generating function agrees with NIST DLMF 18.23.7 after the
original parameter rescaling; it is derived here from the supplied
recurrence, not substituted for it.

List the original roots
$\omega_\nu=(\lambda_{a,b}-c)/i$ in lexicographic order, repeating each
exactly $e$ times, and put
\[
 P(y)=i^{-q}\chi(c+iy)=\prod_{\nu=1}^q(y-\omega_\nu),\qquad
 \Phi:[f(y)]\mapsto[f((S-c)/i)].
 \tag{AKS17}
\]
The inverse is substitution $S=c+iy$, and
$P((S-c)/i)=i^{-q}\chi(S)$. This is an exact algebra isomorphism, including
the complete confluent orders; there is no monomial comparison divisor.
Let $C_s$ have columns $[u_0],\ldots,[u_{q-1}]$. Its nonzero determinant
is $i^{-q(q-1)/2}\prod_{d<q}(h_d^{(s)})^{-1/2}$. Define
\[
 v_j=C_s^{-1}[u_{q+j}],\qquad
 \mathsf K_j=I_q+\sum_{a=0}^{j-1}v_av_a^*,\qquad 0\le j\le q+1.
 \tag{AKS18}
\]
Parseval in the low-degree original orthonormal polynomial basis proves
\[
 \|v_j\|_2^2=
 \frac{\|\operatorname{rem}_P p_{q+j}\|_{m_{0,s}}^2}
      {M_s(q+j)!(b_s)_{q+j}}.
 \tag{AKS19}
\]
This is the full norm of the actual remainder, not a diagonal entry of
an approximated quotient Gram.

\subsection{Newton interpolation with all roots and their orders}
For $0\le a<q$ define $N_a(y)=\prod_{\nu=1}^a(y-\omega_\nu)$, with
$N_0=1$. For every complex $v$, the unique entire-function remainder is
\[
 \operatorname{rem}_P(e^{vy})
   =\sum_{a=0}^{q-1}[\omega_1,\ldots,\omega_{a+1}]e^{v\cdot}\ N_a(y),
 \qquad
 \left|[\omega_1,\ldots,\omega_{a+1}]e^{v\cdot}\right|
       \le\frac{|v|^a}{a!}e^{R|v|}.
 \tag{AKS20}
\]
Here divided differences at repeated nodes mean Hermite divided
differences, not evaluation with a zero denominator.

A proof covering the confluent case is useful. On the standard
$a$-simplex with barycentric coordinates $\theta_1,\ldots,\theta_{a+1}$,
\[
 [\omega_1,\ldots,\omega_{a+1}]e^{v\cdot}
       =v^a\int_{\Delta_a}e^{v\sum_\nu\theta_\nu\omega_\nu}\,d\theta.
\]
The integral of $\prod\theta_\nu^{n_\nu}$ is
$\prod n_\nu!/(a+\sum n_\nu)!$, as follows by iterated integration of
integer beta integrals. Expanding the exponential therefore gives
$\sum_{n\ge0}v^{a+n}h_n(\omega_1,\ldots,\omega_{a+1})/(a+n)!$.
For monomials, the divided difference of $y^{a+n}$ is exactly $h_n$:
this follows by multiplying the geometric series
$\prod(1-\omega_\nu z)^{-1}$ and the Newton division recurrence.
It is a polynomial identity in the nodes and so also holds at repeated
nodes. This proves the integral formula. Iterated Newton division for
polynomials, followed by the convergent exponential series, proves the
first equality in AKS20, including every required derivative.
Finally the simplex has volume $1/a!$ and
$|\sum\theta_\nu\omega_\nu|\le R$, proving its stated bound.

For any $0<t<\pi/2$, put $\tau=2/t$. Positivity of every term of the
cosh series in AKS15 gives
\[
 \|y^a\|_{m_{0,s}}
 \le\sqrt{M_s}(\cos t)^{-b_s/2}\tau^a a!.
\]
Indeed its square is at most
$M_s(\cos t)^{-b_s}(2a)!t^{-2a}$ and
$(2a)!\le4^a(a!)^2$. Expanding the \emph{whole} Newton polynomial, then
using $|e_j(\omega_1,\ldots,\omega_a)|\le\binom ajR^j$, proves
\[
\begin{split}
 \|N_a\|_{m_{0,s}}
 &\le\sqrt{M_s}(\cos t)^{-b_s/2}
        \sum_{j=0}^a\binom ajR^{a-j}\tau^j j!\\
 &=\sqrt{M_s}(\cos t)^{-b_s/2}\tau^a a!
           \sum_{n=0}^a\frac{(R/\tau)^n}{n!}\\
 &\le\sqrt{M_s}(\cos t)^{-b_s/2}\tau^a a!e^{R/\tau}.
\end{split}\tag{AKS21}
\]
There are no omitted polynomial coefficients. Combining AKS20--21 gives
for every complex $v$
\[
 \|\operatorname{rem}_P(e^{vy})\|_{m_{0,s}}
 \le\sqrt{M_s}(\cos t)^{-b_s/2}e^{R(|v|+t/2)}
                       \sum_{a=0}^{q-1}(2|v|/t)^a.
 \tag{AKS22}
\]

\subsection{Cauchy extraction in the actual Gamma generating function}
Fix $0<z<1$ and write $a_z=\operatorname{arctanh}z$.
For $|w|=z$ the defining power series gives
$|\arctan w|\le a_z$ and $|1+w^2|\ge1-z^2$.
Apply the literal finite-dimensional remainder map to AKS16 and use
Cauchy's formula for its $d$th coefficient. The norm triangle inequality
on that circle and AKS22 give
\[
\begin{split}
 \frac{\|\operatorname{rem}_P p_d\|_{m_{0,s}}}{d!}
 \le{}&z^{-d}(1-z^2)^{-b_s/2}\sqrt{M_s}(\cos t)^{-b_s/2}\\
 &\mathrel{}\times e^{R(a_z+t/2)}
                    \sum_{a=0}^{q-1}(2a_z/t)^a.
\end{split}\tag{AKS23}
\]
Cauchy's formula is legitimate for the finite remainder vector: its
entries are holomorphic for $|w|<1$ by the actual matrix generating
identity OPG7. The source norm is taken only after this finite-vector
identity; there is no unjustified integral of a nonintegrable generating
function on the large circle.

Substitution of AKS23 into AKS19 retains the factorial denominator and
cancels the two occurrences of the original mass, giving
\[
 \boxed{\ 
 \|v_j\|_2^2\le\frac{d!}{(b_s)_d}z^{-2d}
 [\cos t(1-z^2)]^{-b_s}e^{2R(a_z+t/2)}
 \left[\sum_{a=0}^{q-1}(2a_z/t)^a\right]^2,
 \quad d=q+j.\ }
 \tag{AKS24}
\]
Adding the \emph{complete} rank-one positive matrices in AKS18 now proves
$\mathsf K_j\preceq \mathcal Z_{s,j}I_q$ after setting $z=4/5$ and $t=t_q$.
Indeed $a_{4/5}=\log3$, $1-z^2=9/25$, and $\cos t_q=\sin(1/q)$, so
the exact scalar obtained is AKS3. The high index $j=q+1$ includes the
original polynomial of degree $2q$.

There is also a root-by-root finite refinement. Put
$\mu_{s,d}=\int y^{2d}dm_{0,s}$,
$R_a=\max_{\nu\le a+1}|\omega_\nu|$, and
\[
\begin{gathered}
 B_{s,a}^{\rm root}=\sum_{d=0}^a
     |e_{a-d}(\omega_1,\ldots,\omega_a)|\sqrt{\mu_{s,d}},\\
 J_s^{\rm root}(z)=\frac{(1-z^2)^{-b_s}}{M_s}
 \left[\sum_{a=0}^{q-1}\frac{a_z^a}{a!}e^{a_zR_a}
                      B_{s,a}^{\rm root}\right]^2.
\end{gathered}\tag{AKS25}
\]
Replace the scalar multiplying
$d!z^{-2d}/(b_s)_d$ in AKS24 by $J_s^{\rm root}(z)$ to get a valid
sharper bound, by the same proof before AKS21 is relaxed. Every moment
is exactly $M_s(2d)![w^{2d}](\cos w)^{-b_s}$. Thus this refinement keeps
every actual root coefficient and a finite original moment formula.
AKS21 bounds it by the explicit scalar in AKS24. The proof never replaces
the source determinant by a product of diagonal moments: AKS24 is a
full-vector operator bound, and the next step restricts the inverse of
the full covariance.

\section{The four original kernel Grams and the evaluated scale}
\label{sec:AKS-volume}

The original orthonormal source map at degree $q+j-1$ is
$C_s[I_q,v_0,\ldots,v_{j-1}]$. Its minimum lift of $C_sx$ is
$[I_q,v_0,\ldots,v_{j-1}]^*\mathsf K_j^{-1}x$. It has that image and is
orthogonal to the full relation kernel, so Pythagoras proves
\[
 G_{q+j-1}^{(s)}=C_s^{-*}\mathsf K_j^{-1}C_s^{-1},\qquad
 H_j:=H_{K,q+j-1}^{(s)}=T_s^*\mathsf K_j^{-1}T_s,\quad T_s=C_s^{-1}I.
 \tag{AKS26}
\]
This verifies the source-minimum step in the coefficient determinant
AKS13--14. Neither a cross block nor any original relation is removed.

Since $I_q\preceq\mathsf K_j\preceq \mathcal Z_{s,j}I_q$, inversion and compression
give $\mathcal Z_{s,j}^{-1}H_0\preceq H_j\preceq H_0$, where $H_0=T_s^*T_s$.
Define $d_j=-\log\det(H_0^{-1/2}H_jH_0^{-1/2})$. Congruence and all
$r$ eigenvalues give
\[
 0=d_0\le d_1\le\cdots\le d_{q+1},\qquad
 0\le d_j\le r\mathcal L_{s,j},\qquad
 F_K^{(s)}=d_q+d_{q+1}-d_1.
 \tag{AKS27}
\]
The identity uses precisely the four endpoints in AKS2.
In particular it proves AKS4, as well as the sharper interval
\[
 \boxed{\quad d_1\le F_K^{(s)}\le
             r(\mathcal L_{s,q}+\mathcal L_{s,q+1})-d_1.\quad}
 \tag{AKS28}
\]
For a directly calculable expression for its lower endpoint set
$V=\|v_0\|^2$ and
$\theta=v_0^*T_s(T_s^*T_s)^{-1}T_s^*v_0$. Then
\[
 0\le\theta\le V,\qquad
 d_1=\log\frac{1+V}{1+V-\theta}.
 \tag{AKS29}
\]
Indeed the inverse of $I+v_0v_0^*$ is
$I-v_0v_0^*/(1+V)$; take its full compressed determinant. For $r=0$
set $\theta=d_1=0$. In GMB/OPR notation $V=t_0$ and $V-\theta=\xi_0$,
so this is exactly the first original section-corrected kernel increment,
not an independently chosen kernel observation.

Here is an evaluated finite remainder for AKS6. Put
\[
\begin{gathered}
 a_\infty=\frac{4\log3}{\pi},\qquad
 c_0=\log\frac{25\log3}{4\pi}>0,\\
 E_{s,q}(R)=\log2+\frac4{\pi-2/q}
  +\frac{s}{2}\log\frac{25\pi q}{18}
  +2R\left(\log3+\frac\pi4\right)\\
 \hspace{23mm}+\log((q+1)(4q+1))-2\log(a_\infty-1).
\end{gathered}\tag{AKS30}
\]
Then for both source orders, every $q\ge36$, and the same actual kernel,
\[
 \boxed{\quad
 \mathcal L_{s,q},\mathcal L_{s,q+1}\le2c_0q+E_{s,q}(R),\qquad
 0\le F_K^{(s)}\le4rc_0q+2rE_{s,q}(R).
 \quad}\tag{AKS31}
\]
To verify every factor, $b_s\ge1/2$ gives
\[
 \frac{d!}{(b_s)_d}\le\frac{4^d}{\binom{2d}{d}}\le2d+1.
\]
The last inequality uses that the central binomial is the largest of the
$2d+1$ positive coefficients whose sum is $4^d$. Thus the sum in AKS3
is at most $(q+1)(4q+1)(5/4)^{4q}$. Also
$\mathfrak s_q\le \mathfrak a_q^q/(a_\infty-1)$ and
\[
 2q\log \mathfrak a_q\le2q\log a_\infty+\frac4{\pi-2/q},\qquad
 \sin(1/q)\ge\frac2{\pi q}.
\]
The first follows from $-\log(1-x)\le x/(1-x)$ at $x=2/(\pi q)$;
the second is the chord bound for sine on $[0,\pi/2]$.
Use $t_q\le\pi/2$ in the remaining exponential. The order-$q$ terms
combine to
$4q\log(5/4)+2q\log a_\infty=2c_0q$.
The upper bound for the logarithm of the added positive scalar is
positive, so $\log(1+e^x)\le\log2+\max(0,x)$ proves the first bound in
AKS31. AKS28 proves the second.

On the original simple family $q=(k+1)^2$, $R=O_h(k)$ and
$r=8k-16$. Consequently $E_{1,q}(R)=O_h(k+\log q)$ and
$E_{k,q}(R)=O_h(k\log q)$. Dividing AKS31 by $kq$ proves AKS6, with
$32c_0=\kappa_*$. In fact the fixed-source upper bound is
\[
 F_K^{(1)}\le(32k-64)c_0q+(16k-32)E_{1,q}(R),
 \tag{AKS32}
\]
whose explicit second term is $O_h(k^2)$. The tensor-source version has
an $O_h(k^2\log q)$ upper remainder. These are remainders of upper
bounds, not assertions about the actual signed error or leading value.
The proof is uniform in all the period parameters in the stated domain:
no norm or condition number of a period matrix occurs in AKS30.

\paragraph{The precise earlier loss.}
OKV15 had already retained
$\|v_j\|^2\le q^22^{6q}(2D_*q)^{2d}/(d!(b_s)_d)$, $d=q+j$.
Replacing its denominator by $1/2$ caused the old $q\log q$ exponent.
Even without AKS20--24 one obtains the valid repair
\[
 \|\mathsf K_{q+1}\|\le
 1+(q+1)q^22^{6q}\frac{(4D_*q)^{4q}}{(4q)!},\qquad D_*\ge3.
 \tag{AKS33}
\]
Here $d!(b_s)_d\ge(2d)!/4^d$ and the remaining factorial ratio increases
up to $d=2q$ for $D_*\ge3$. The elementary inequality
$\log n!\ge n\log n-n+1$ already makes the logarithm of AKS33 $O_h(q)$.
The Newton calculation strengthens this repair to the packet-independent
limsup constant AKS5, while retaining $R$ in its finite remainder.


\paragraph{Exact enclosure of the displayed numerical ceiling.}
For $x>1$ and $z=(x-1)/(x+1)$ the integrated geometric series gives
\[
 0\le \log x-2\sum_{j=0}^{39}\frac{z^{2j+1}}{2j+1}
      \le\frac{2z^{81}}{81(1-z^2)}.
\]
For $0<x<1$, the alternating sum through $j=39$ for $\arctan x$
has error between zero and $x^{81}/81$. The tangent addition formula,
with the angle in $(0,\pi/2)$, proves
$\pi=16\arctan(1/5)-4\arctan(1/239)$: the tangent of
$4\arctan(1/5)-\arctan(1/239)$ is exactly one. Use the resulting
rational lower and upper bounds for $\pi$ and for $\log3$, then apply
the first displayed bound to the two endpoints of
$25\log3/(4\pi)$. Multiplication by $32$ gives, with outward decimal
rounding, the interval
\[
 \kappa_*\in[25.0207809764994871,\ 25.0207809764994872].
\]
The accompanying rational-arithmetic implementation checks these
inequalities by integer cross multiplication, rather than a floating-point
logarithm. This proves the less finely rounded strict interval AKS5.

\section{The complete arithmetic and mixed receivers}
\label{sec:AKS-transfer}

\subsection{The unchanged finite arithmetic allowances}
Use exactly AMT4's constants. To prevent a loss hidden in order notation,
write them here:
\[
\begin{gathered}
 \alpha=\pi/2,\ p=8m-9/2,\ B_h=42+8m,\ M=21+4m,\\
 \vartheta_h=\int_{-1}^1w_h(y)dy,\quad
 M_\alpha=\int e^{\alpha y}w_h(y)dy,\quad
 w_h(y)=|v_h(1/2+iy)|^2/(2\pi),\\
 c_\Gamma=\sqrt2e^{-7/6},\quad C_\Gamma=\sqrt{2\sqrt5}e^{2/3},\\
 a_k=\frac{c_h\vartheta_h^{k-3}e^{-\alpha(k-3)}(k-2)^{-B_h}}
                    {C_\Gamma2^{B_h/2}},\quad
 A_k=\frac{C_h^{\rm tilt}}{c_\Gamma}k^{p+1}M_\alpha^{k-1},\\
 D_N=[1+4(N+M)^2]^M,\quad L_N^{\rm ar}=\log(A_kD_N/a_k),\\
 E_k^+=L_{q-1}^{\rm ar}+L_q^{\rm ar},\qquad
 E_k^-=L_{2q-1}^{\rm ar}+L_{2q}^{\rm ar}.
\end{gathered}\tag{AKS34}
\]
The constants $c_h,C_h^{\rm tilt}$ are the original TW/BT envelope
constants, not new choices. Their actual inequalities and complete
source construction are preserved in AMT4--6 and the current source
closure. Minimize AMT6 on the same remainder fibre and then restrict to
$I(K)$. AMT9--20 gives, with the original arithmetic source $w_h^{*k}$,
\[
 -rE_k^+\le F_K^{\rm ar}-F_K^{(1)}\le rE_k^-.
\]
Let $\underline U_s=d_1^{(s)}$ and
$\overline U_s=r(\mathcal L_{s,q}+\mathcal L_{s,q+1})-d_1^{(s)}$. The new finite conclusion is
\[
 \boxed{\quad
 \max\{0,\underline U_1-rE_k^+\}\le F_K^{\rm ar}
                  \le\overline U_1+rE_k^-.
 \quad}\tag{AKS35}
\]
This improves AMT37's absolute upper bound. For the sharper correlated
version retain $d_{E,N}=q\log A_k-
\log(\det G_N^{\rm ar}/\det G_N^{(1)})$ and
$c_{X,N}=\min\{r_XL_N^{\rm ar},d_{E,N}\}$, where $r_K=r$, $r_B=q-r$.
Put $C_X^-=c_{X,q-1}+c_{X,q}$,
$C_X^+=c_{X,2q-1}+c_{X,2q}$. AMT18 then gives
\[
\begin{split}
 \max\{0,\underline U_1-C_K^-,
                  \underline U_1+\delta_k^\sigma-C_B^+\}
 &\le F_K^{\rm ar}\\
 &\le\overline U_1+
            \min\{C_K^+,\delta_k^\sigma+C_B^-\}.
\end{split}\tag{AKS36}
\]
One may replace $d_{E,N}$ in the caps by the original full-source
$\Xi_{k,N}$ or its proved finite majorant. The scalar
$\delta_k^\sigma=\mathcal B_k^{\rm ar}-\mathcal B_k^\sigma$ retains
exactly AMT21:
\[
\begin{gathered}
 \max\{-qE_k^+,-\Xi_{k,q-1}-\Xi_{k,q}\}
       \le\delta_k^\sigma\\
       \le\min\{qE_k^-,\Xi_{k,2q-1}+\Xi_{k,2q}\}.
\end{gathered}\tag{AKS37}
\]
Since $E_k^\pm=O_h(k+\log q)$ by the exact AMT22 expression, AKS35 and
AKS32 prove the arithmetic version of AKS6, with an explicit
$O_h(k^2)$ upper remainder. Nonnegativity follows from the actual
nested source minima, as in AMT15, not from comparison with a possibly
negative lower endpoint.

\subsection{The original Gamma multiplier, not a replacement higher source}
Retain $t_j=2j+1/2$, $f_l(S)=\prod_{j=0}^{l-1}(S-c-t_j)$ and
\[
 dm_{0,k}=\beta_k|f_l(c+iy)|^2d\sigma,\quad
 \beta_k=\frac{(2\pi)^{k/2}}{\sqrt2\Gamma(k/2)},\quad
 T_l=\prod_{j=0}^{l-1}t_j^2,\quad
 U_{l,N}=\prod_{j=0}^{l-1}[4(N+j+1)^2+t_j^2].
 \tag{AKS38}
\]
The exact source map is $\mathcal P_N\to f_l\mathcal P_N\subset
\mathcal P_{N+l}$, $P\mapsto f_lP$, with inverse division on this
specified subspace. Let $\mathsf E_\chi$ be the complete divided-jet
map and $U_f$ its Leibniz multiplication matrix,
$(U_f)_{r,a}=f_l^{(r-a)}(\lambda)/(r-a)!$ for $r\ge a$.
AMT25--28 gives the precise transported observation and kernel
\[
\begin{gathered}
 \widetilde\Lambda=\Lambda\mathsf E_\chi^{-1}U_f^{-1},\quad
 \widetilde I=U_f\mathsf E_\chi I,\quad
 \ker\widetilde\Lambda=U_f\mathsf E_\chi K,\\
 C_N=\mathscr K_{\chi,N+l}
   -\mathscr K_{\chi,f;N+l}\mathscr K_{f,N+l}^{-1}
                               \mathscr K_{f,\chi;N+l},\\
 H_{K,N}^{(k)}=\beta_k\widetilde I^*C_N^{-1}\widetilde I.
\end{gathered}\tag{AKS39}
\]
Each covariance is the original finite sum of full polynomial jet outer
products divided by its original monic norm. These formulas retain the
extra $f$ constraints and their full cross covariance. In particular
$M_fK=K$ is not assumed. AKS26 calculates the same tensor metric directly
in its own orthonormal source; the two representations agree by their
identical source norm and identical constrained minimum.

Define
\[
 w_N=\log(U_{l,N}/T_l),\quad
 \Omega^{\rm low}=w_{q-1}+w_q,\quad
 \Omega^{\rm high}=w_{2q-1}+w_{2q}.
\]
The already proved same-coefficient comparison AMT42--44 retains all
mass factors before cancellation and gives
\[
 -r\Omega^{\rm low}\le\Delta_K^\Gamma:=F_K^{(1)}-F_K^{(k)}
                         \le r\Omega^{\rm high}.
 \tag{AKS40}
\]
It therefore sharpens the two independent absolute bounds by
\[
\begin{split}
 \max\{\underline U_k,\underline U_1-r\Omega^{\rm high},0\}
       &\le F_K^{(k)}\\
       &\le\min\{\overline U_k,\overline U_1+r\Omega^{\rm low}\}.
\end{split}\tag{AKS41}
\]
The original finite estimate
$0\le w_N\le l\log[1+16(N+l)^2]$ proves that AKS40 is $o(kq)$ on the
actual rank-$r$ domain. Along with AKS35 this preserves the completed
fact that all three source kernels differ by $o(kq)$. AKS4, not AKS40,
is the new absolute-scale calculation.

\subsection{Boundary, QGQ and arithmetic receiving identities}
At each endpoint retain the original minimum section
$S_N^{(s)}=(G_N^{(s)})^{-1}\Lambda^*
(\Lambda(G_N^{(s)})^{-1}\Lambda^*)^{-1}$. The full Schur identity is
\[
 \det(I^*G_N^{(s)}I)\det Q_N^{(s)}
       =|\det[I,S_*]|^2\det G_N^{(s)},\qquad
 Q_N^{(s)}=(\Lambda(G_N^{(s)})^{-1}\Lambda^*)^{-1}.
 \tag{AKS42}
\]
Thus the new finite absolute boundary interval is
\[
 \max\{0,\mathcal B_k^{(s)}-\overline U_s\}
       \le F_B^{(s)}\le\mathcal B_k^{(s)}-\underline U_s.
 \tag{AKS43}
\]
The baselines in this formula are the \emph{complete} original BSL/HBL
objects, with their full finite relation errors. No $q^2$ coefficient is
assigned to $K$; its proved coefficient there remains zero.

Write $\mathcal H_k=\mathcal B_k^\sigma-\mathcal B_k^0=-\mathfrak T_k$.
The exact full identities still read
\[
\begin{gathered}
 F_B^{(1)}-F_B^{(k)}=\mathcal H_k-\Delta_K^\Gamma,\\
 \boxed{\quad\mathcal B_k^{\rm ar}
   =F_K^{(k)}+F_B^{(k)}+\mathcal H_k+\delta_k^\sigma
   =F_K^{(1)}+F_B^{(1)}+\delta_k^\sigma.\quad}
\end{gathered}\tag{AKS44}
\]
The Gamma interval is not re-estimated in this argument. For either
complete QRI1 or QRI2 interval $[\omega^-,\omega^+]$ for its already
calculated centre $W_k$, retain its exact source errors
$e_0\in[-a_0,b_0]$ and $e_q+e_{q+1}\in[-b_k^+,b_k^-]$. Then
\[
 \underline{\mathcal H}_k=-\omega^+-b_k^--a_0,
 \qquad\overline{\mathcal H}_k=-\omega^-+b_k^++b_0,
 \qquad \mathcal H_k\in
       [\underline{\mathcal H}_k,\overline{\mathcal H}_k].
 \tag{AKS45}
\]
These are exactly QRI3, not new unspecified error choices. The complete
QRI1--9, QGT and QGQ definitions remain verbatim in each updated receiver;
the new interval is intersected with all their retained tighter finite
intervals. In particular the simple-quartet nonzero error radius stays
\[
 \limsup\left|
 \frac{\mathcal H_k+lq\mathcal C_\Gamma}{q}
       +\frac{\partial_\alpha\mathcal L(2,\pi)}8
       -\frac{\log2}{4}\right|\le\frac{A_{\rm QGT}}{\sqrt2}.
 \tag{AKS46}
\]
No zero-radius or $o(q)$ arithmetic assertion is introduced.

For example the complete finite absolute defect relative to the tensor
boundary, with the actual Gamma and arithmetic scalars retained, is
\[
 \underline U_k+\mathcal H_k+\delta_k^\sigma
 \le\mathcal B_k^{\rm ar}-F_B^{(k)}
 \le\overline U_k+\mathcal H_k+\delta_k^\sigma.
 \tag{AKS47}
\]
Inserting AKS41 makes this sharper. Replacing $\mathcal H_k$ by the
appropriate endpoints of AKS45 and $\delta_k^\sigma$ by AKS37 gives its
fully inherited finite scalar enclosure with the same signs.
Also $F_B^{\rm ar}=\mathcal B_k^\sigma+\delta_k^\sigma-F_K^{\rm ar}$;
subtract the interval AKS35 or AKS36 to get its new absolute boundary
interval. This retains the exact correlations rather than taking
independent hypothetical source errors.

The existing limit is
$\mathcal H_k/(kq)\to-\mathcal C_\Gamma/4$ and
$\delta_k^\sigma/(kq)\to0$. Therefore the new interval implies
\[
\begin{gathered}
 \liminf_k\inf_{|u|\ge R_*}
 \frac{\mathcal B_k^{\rm ar}-F_B^{(k)}}{kq}
          \ge-\mathcal C_\Gamma/4,\\
 \limsup_k\sup_{|u|\ge R_*}
 \frac{\mathcal B_k^{\rm ar}-F_B^{(k)}}{kq}
          \le\kappa_*-\mathcal C_\Gamma/4.
\end{gathered}\tag{AKS48}
\]
All original branches are included in the infima and suprema. This does
not decide its sign. The already proved boundary \emph{difference}
limits of AMT49 are unchanged. Every subsequential limit of the absolute
normalized kernel lies in $[0,\kappa_*]$ and is shared by its three source
versions when they are evaluated on the same sequence of period
parameters. Existence of a single limit is not proved by boundedness.

Finally the HC residual remains exactly
$\mathcal R_k=\mathcal B_k^{\rm ar}-4q\log(D_hk)$, with its original
$D_h$ and original nonnegative contraction, phase and trace terms.
AKS44 changes none of them. The supplied positive $q^2$ baseline and
its complete Gamma correction are not re-presented as new results here.

\section{Actions, support faces, and the actual arithmetic class}
\label{sec:AKS-actions}

The coefficient formula AKS12 supplies all four blocks of the same sum
action: $P_KYP_K$, $P_KY(I-P_K)$, $(I-P_K)YP_K$ and
$(I-P_K)Y(I-P_K)$. There is no claim that either off-diagonal block
vanishes. In the actual minimum-section coordinates the corresponding
formula remains
\[
 [I,S_N]^{-1}Y[I,S_N]=
 \begin{pmatrix}
 H_{K,N}^{-1}I^*G_NYI&H_{K,N}^{-1}I^*G_NYS_N\\
 \Lambda YI&\Lambda YS_N
 \end{pmatrix}.
 \tag{AKS49}
\]
Its proof is multiplication by the inverse coordinate maps
$H_{K,N}^{-1}I^*G_N$ and $\Lambda$. In the invariant-ring model, the same
action is the derivation
$\mathcal Df=(\mathsf H\Omega\mathsf H^{-1}x)\cdot\nabla f$,
$\Omega=\operatorname{diag}(\gamma-i\delta,-\gamma-i\delta,
\gamma+i\delta,-\gamma+i\delta)$, with $\mathcal DQ=0$.
For the actual conductor $\mathcal A_Q=\prod_{a=1}^4Q_a$, its full defect
is $\mathcal D\mathcal A_Q=\sum_{a=1}^4(\mathcal DQ_a)
\prod_{b\ne a,\,1\le b\le4}Q_b$, exactly OQC22. Thus this coefficient
realization preserves the original action leaving the observed dual.

In particular the section-corrected vector is still OPR2's
\[
 C_s r_j^{\rm ker}=e_{q+j}-S_{q+j-1}^{(s)}\Lambda e_{q+j}
   =I\bigl(\zeta_{q+j}-\kappa S_{q+j-1}^{(s)}\Lambda e_{q+j}\bigr),
\]
with its full denominator $1+\xi_j$ in the kernel increment. The operator
proof estimates the complete covariance and its actual restriction, so
none of these coupled terms is replaced by a raw fixed-section vector.
The all-iterate kernel stays
$\bigcap_{a=0}^{q-1}\ker(\Lambda Y^a)$; nothing in AKS4 declares it zero.

On the existing arithmetic coefficient realization retain
$\mathcal V_{h,k}P=P(D_1+\cdots+D_k)F_h^{\otimes k}$ and
$J^{(k)}\mathcal V_{h,k}=\eta\operatorname{rem}_\chi$, where
$\eta=U^{\otimes k}\iota$ is the original injection. AKS17 and the
Newton coefficients compute precisely this same remainder. Any two
polynomial lifts differ by $\chi Q$ in the original relation space,
so their original arithmetic cohomology classes agree. The specified
primitive of the difference is EP.48, with the full differential and
cochain signs proved in AKP3--5 below; it depends on the lift difference.
The source comparison to $w_h^{*k}$ is the identity on original
polynomial coefficients, followed by the actual minima as in AMT8;
it is not a replacement arithmetic source.

The scalar object stays $G(R)=\{\tau\}\sqcup R^\bullet$, with supported
zero $0_R^\bullet$, injective amplitude/support pair, and infinite
integer target. Every coefficient map above is applied on each original
present support label and outer leg. It commutes with the specified
coefficient inclusions because both compositions apply the same linear
map to the same coefficient; a present zero is not converted into
absence. No new base object, purity condition, or identification of the
sum action with a Frobenius action is used. The actual distinguished
class and its higher-jet directions remain in the full fibre, even
when the observation kills them.

\section{Exact period domain and the interior alternative}
\label{sec:AKS-radius}

For reference the inherited radius used above is completely specified
as follows; these are AMT33--34/RPC17--23, with no new constants. Put
\[
\begin{gathered}
 \beta=2(\gamma^2-\delta^2),\quad\eta=(\delta^2+\gamma^2)^2,\quad
 a=v_h(\rho)\ne0,\\
 c_3=(a^{-2}-\bar a^{-2})/(4i\delta\gamma),\quad
 c_1=(a^{-2}+\bar a^{-2})/2-(\delta^2-\gamma^2)c_3,\\
 g_j=5^{j/5-1}e^{\pi ij/5}\Gamma(j/5),\quad
 g_+=\max_j|g_j|,\quad g_-^{-1}=\max_j|g_j|^{-1},\quad
 \kappa=\max_{j<4}|g_{j+1}/g_j|,\\
 E_j^{\rm ray}=\frac85\int_0^\infty r^j(r^3/3+r)
       e^{-r^5/5+r^3/3+r}dr,\quad
 C_T=2\left(\sum_{j=0}^3(E_j^{\rm ray})^2\right)^{1/2},\\
 C_K=4C_Tg_-^{-1}+2g_+g_-^{-1}+2g_+g_-^{-2}C_T,\\
 C_{K,3}=C_K(3\kappa^2+3\kappa C_K+C_K^2),\quad
 \epsilon_* =\min\{1,(2C_Tg_-^{-1})^{-1}\},\quad B_* =\beta+\eta.
\end{gathered}\tag{AKS50}
\]
Then
\[
R_*=
\begin{cases}
\displaystyle\left[\max\left\{1,\frac{B_*}{\epsilon_*},
\frac{2(|c_3|C_{K,3}B_*+|c_1|(\kappa+C_K))}{|c_3g_4/g_1|}
\right\}\right]^{5/2},&c_3\ne0,\\[4pt]
\displaystyle\left[\max\left\{1,\frac{B_*}{\epsilon_*},
\frac{2C_KB_*}{|g_2/g_1|}\right\}\right]^{5/2},&c_3=0.
\end{cases}\tag{AKS51}
\]
The actual unit implies $(c_1,c_3)\ne(0,0)$ and all integrals converge.
RPC's full period comparison proves a nonzero $(4,1)$ or $(2,1)$
commutator entry, respectively, for $|u|\ge R_*$ on every branch. It
therefore proves the five-member orbit needed in AKS11.

Inside this radius the same coefficient determinant AKS12--14 and the
operator bound AKS24 remain valid. The orbit is decided by the original
noninvariant quadratic coefficients, not by the radius alone. Since
$Q$ is nondegenerate, its projective orbit has either one or five
members. A one-member orbit must have character one: its character has
both fourth and fifth powers equal to one, by the quadratic determinant
and $D^5=I$. In that case $Q=a'x_1x_4+b'x_2x_3$ with $a'b'\ne0$, and
\[
 b=d_{5,k,0}-d_{5,k-2,0}
   =\left\lfloor\frac{(k+1)^2}{5}\right\rfloor
           +\mathbf1_{k\equiv0\text{ or }3\ (5)},\qquad r=q-b.
 \tag{AKS52}
\]
If a noninvariant coefficient is nonzero, AKS11 applies instead. These
are the complete OQC alternatives. The bound AKS31 always uses that
actual $r$; it does not extend the $8k-16$ specialization to an
unevaluated exceptional point. This calculation does not assert that
the interior exceptional set is empty.

\section{Verification scope and the remaining absolute coefficient}
\label{sec:AKS-scope}

The new mathematical content is AKS20--32 and its application to the
actual coefficient determinant AKS12--14, followed by the finite
receiving improvements AKS35--48. The complete source and orbit
identifications used in that application are the supplied OCS/OQC/OQX,
RPC, OPG/OPR and AMT results; their exact predecessors are preserved.
AKS33 pinpoints a repair obtainable immediately from the earlier proof,
whereas the Newton bound supplies the sharper explicit constant.

The bound \emph{at} scale $kq$ is now finite and uniform. What is not
proved is a matching lower coefficient, a unique limit for
$F_K^{(1)}/(kq)$, or a sign for the interval AKS48. Those questions involve
the orientation of the original period kernel within the inverse
remainder covariance; replacing that compression by a scalar bound
loses precisely that information. This is the finite-bounds outcome
allowed by the current request, not a claim to have evaluated an
asymptotic coefficient. The source differences are already completed
inputs and are not renamed as the remaining problem.

The accompanying exact checks distinguish finite algebraic diagnostics
from the analytic proof for the full family. They do not assert that a
diagnostic root is a zero of $\xi$, that a diagnostic matrix equals an
actual period matrix, or that a new Lean theorem has been checked.
No RH conclusion is drawn from this kernel bound.
